

\cleardoublepage
\thispagestyle{empty}
\bigtitle{S.\@ Ephræm Syri Conf.\@ et Eccl.\@ Doct.}
\addcontentsline{toc}{section}{XVIII Junii} \phantomsection
\duplex

\header{xviii junii 2026}

\office{ad primam}
\addcontentsline{toc}{subsection}{Ad Primam} \phantomsection

\pateravecredo

\gscore[]{℣.}{Deus_in_adiutorium_Tonus_ferialis_no_laus_tibi}

\gscore[Hy.]{3.}{hy_jam_lucis_orto_sidere_in_festo_s._cordis_jesu_solesmes_1961}

\gscore[Ant.]{7 a.}{in_loco_pascuae}
\psalmus{2}{22_7a}{22_7}
\psalmus{71.\textsc{i}}{71_1_7a}{71_1_7}
\psalmus{71.\textsc{ii}}{71_2_7a}{71_2_7}
\smallscore{an_in_loco_pascuae_ibi_dominus_solesmes_1961}

\capitulum{1 Tim. 1, 17}
\lettrine{R}{e}gi sæculórum immortáli et invisíbili,~† soli Deo honor et glória~* in sǽcula
sæculórum.

\rr Deo grátias.

\gscore[]{Resp. brev.}{rb_christe_fili_dei_vivi_sc}

\smallscore{versiculus_ad_primam}

  \vv Dóminus vobíscum. \rr Et cum spíritu tuo.
  
  \oratio
  
  Orémus.
  
 \lettrine{D}{o}mine Deus omnípotens, qui ad princípium huius diéi nos perveníre
fecísti: tua nos hódie salva virtúte; ut in hac die ad nullum declinémus peccátum,
sed semper ad tuam justítiam faciéndam nostra procédant elóquia, dirigántur
cogitatiónes et ópera. Per Dóminum.

\rr Amen.

\vv Dóminus vobíscum. \rr Et cum spíritu tuo.

\smallscore{BD_ad_completorium}

\smalltitle{Martyrologium.}

\begin{enpars}
\lettrine[lraise=0.2]{J}{u}ne 19\textsuperscript{th}, 2026, the fifth day of the Moon, were born into the better life:
The Octave of the most Sacred Heart of Jesus.
At Florence, St. Juliana Falconieri, virgin, foundress of the Sisters of the Order of the Servants of the Blessed Virgin Mary, who was placed among the holy virgins by the Sovereign Pontiff, Clement \scspace{xii}.
At Milan, (in the first century,) the holy martyrs Gervase and Protase, brothers. The former, by order of the judge Astasius, was scourged with leaded whips for so long that he expired. The latter, after being scourged with rods, was beheaded. Through divine revelation their bodies were found by St. Ambrose. They were partly covered with blood, and as free from corruption as if they had been put to death that very day. When the translation took place, a blind man recovered his sight by touching their relics, and many persons possessed by demons were delivered.
At the monastery in the valley of Castro in Piceno, the birthday of St. Romuald, anchoret, a native of Ravenna. He was the founder of the Camaldolese monks, and he restored and greatly extended monastic discipline, which was much relaxed in Italy. His feast is observed on the 7\textsuperscript{th} of February, on which day his sacred relics were transferred to Fabriano.
At Arezzo in Tuscany, the holy martyrs Gaudentius, bishop, and Culmatius, deacon, who were murdered by the furious heathen, during the reign of Valentinian.
Also, St. Boniface, martyr, a disciple of blessed Romuald, who was sent by the Roman Pontiff, Gregory \textsc{v}, to preach the Gospel in Russia. Having passed through fire uninjured, and baptized the king and his people, he was killed by the enraged brother of the king, and thus gained the palm of martyrdom which he ardently desired.
At Ravenna, (likewise in the first century,) St. Ursicinus, martyr, who remained constant through many torments in the confession of martyrdom by being beheaded.
At Sozopolis, under the governor Domitian, during the persecution of Trajan, St. Zosimus, martyr, who suffered bitter tortures, was beheaded, and thus triumphantly went to heaven.
\end{enpars}

\vv Et álibi aliórum plurimórum san\-ctórum Mártyrum et Confessórum, atque san\-ctárum Vírginum.

\rr Deo grátias.

\vv Pretiósa in conspéctu Dómini. \rr Mors Sanctórum ejus.

Sancta María et omnes Sancti intercédant pro nobis ad Dóminum, ut nos
mereámur ab eo adjuvári et salvári, qui vivit et regnat in sǽcula sæculórum.

\rr Amen.

\vv Deus in adjutórium meum inténde.

\rr Dómine, ad adjuvándum me festína.

\vv Deus in adjutórium meum inténde.

\rr Dómine, ad adjuvándum me festína.

\vv Deus in adjutórium meum inténde.

\rr Dómine, ad adjuvándum me festína.

\vv Glória Patri, et Fílio,~* et Spirítui Sancto.

\rr Sicut erat in princípio, et nunc, et semper,~* et in sǽcula sæculórum. Amen.

\vv Kýrie eléison.

\rr Christe eléison. Kýrie eléison.

\vv Pater noster \textit{secreto.}

\vv Et ne nos indúcas in tentatiónem:

\rr Sed líbera nos a malo.

\vv Réspice in servos tuos, Dómine, et in ópera tua, et dírige fílios eórum.

\rr Et sit splendor Dómini Dei nostri super nos, et ópera mánuum nostrárum
dírige super nos, et opus mánuum nostrárum dírige.

\vv Glória Patri, et Fílio,~* et Spirítui Sancto.

\rr Sicut erat in princípio, et nunc, et semper,~* et in sǽcula sæculórum. Amen.

Orémus.

\lettrine{D}{i}rígere et sanctificáre, régere et gubernáre dignáre, Dómine Deus, Rex
cæli et terræ, hódie corda et córpora nostra, sensus, sermónes et a\-ctus nostros in
lege tua, et in opéribus mandatórum tuórum: ut hic et in ætérnum, te auxiliánte,
salvi et líberi esse mereámur, Salvátor mundi: Qui vivis et regnas in sǽcula
sæculórum.

\rr Amen.

\vv Jube, domne, benedícere.

\smalltitle{Benedictio.}

Dies et actus nostros in sua pace dispónat Dóminus omnípotens.

\rr Amen.

\lectiobrevis{Sap. 10, 10}
\lettrine{J}{u}stum dedúxit Dóminus per vias re\-ctas, et osténdit illi regnum Dei, et dedit illi sciéntiam san\-ctórum: honestávit illum in labóribus, et complévit labóres illíus.

\vv Tu autem, Dómine, miserére nobis.

\rr Deo grátias.

\vv Adjutórium nostrum in nómine Dómini.

\rr Qui fecit cælum et terram.

\vv Benedícite.

\rr Deus.

\lettrine{D}{o}minus nos benedícat, \cc{} et ab omni malo deféndat, et ad vitam perdúcat ætérnam. Et fidélium ánimæ per misericórdiam Dei requiéscant in pace.

\rr Amen.

\office{ad missam}
\addcontentsline{toc}{subsection}{Ad Missam} \phantomsection
\proper{Introitus}{Eccli. 15, 5; Ps. 91, 2}

\texteparacol{\lettrine{I}{n}  médio Ecclésiæ apéruit os ejus: et implévit eum Dóminus spíritu sapiéntiæ et intelléctus: stolam glóriæ índuit eum.
℣. Bonum est confitéri Dómino: et psállere nómini tuo, Altíssime.
Glória Patri. Sicut erat.
In médio.}{english}{\noindent In the midst of the assembly he opened his mouth; and the Lord filled him with the spirit of wisdom and understanding; he clothed him with a robe of glory. ℣. It is good to give thanks to the Lord, to sing praise to thy name, Most High. Glory be to the Father. As it was in the beginning. In the midst.}

\masspart{Oratio}

\texteparacol{\lettrine{D}{e}us, qui Ecclésiam tuam beáti E\-phræm Confessóris tui et Do\-ctóris mira eruditióne et præcláris vitæ méritis illustráre voluísti:~† te súpplices exorámus; ut i\-pso intercedénte,~* eam advérsus erróris et pravitátis insídias perénni tua virtúte deféndas.
Per Dóminum.

\lettrine[ante=2]{D}{e}us, qui nobis in Corde Fílii tui, nostris vulneráto peccátis, infinítos dile\-ctiónis thesáuros misericórditer largíri dignáris:~† concéde quǽsumus; ut illi devótum pietátis nostræ præstántes obséquium,~* dignæ quoque satisfa\-ctiónis exhibeámus offícium.

\lettrine[ante=3]{P}{r}æsta, quǽsumus omnípotens Deus:~† ut qui san\-ctórum Mártyrum tuórum Marci et Marcelláini natalítia cólimus,~* a cun\-ctis malis imminéntibus eórum intercessiónibus liberémur. Per Dóminum.}{english}{\noindent  O God, who willed to enlighten thy church by the wondrous learning and glorious merits of the life of blessed Ephrem, thy confessor and Doctor, we humbly pray You, that, by his pleading, thou will shield her with thy lasting power against the snares of error and evil. Through our Lord.

\noindent 2. O God, who hast suffered the heart of thy Son to be wounded by our sins, and in that very heart hast bestowed on us the abundant riches of thy love: grant that the devout homage of our hearts, which we render unto him, may of thy mercy be deemed a recompence acceptable in thy sight.

\noindent 3. Grant, we beseech thee, almighty God, that we who celebrate the anniversary of the death of thy holy martyrs, Mark and Marcellianus, may by their intercession be delivered from all the evils that threaten us. Through our Lord.}

\reading{Epistola}{2 Tim. 4, 1-- 8}

\texteparacol{\lettrine{C}{a}ríssime: Testíficor coram Deo, et Jesu Christo, qui judicatúrus est vivos et mórtuos, per advéntum i\-psíus et regnum ejus: prǽdica verbum, insta opportúne, importune: árgue, óbsecra, íncrepa in o\-mni patiéntia, et doctrína. Erit enim tempus, cum sanam doctrínam non sustinébunt, sed ad sua desidéria, coacervábunt sibi magistros, pruriéntes áuribus, et a veritáte quidem audítum avértent, ad fábulas autem converténtur. Tu vero vígila, in ó\-mnibus labóra, opus fac Evangelístæ, ministérium tuum imple. Sóbrius esto. Ego enim jam delíbor, et tempus resolutiónis meæ instat. Bonum certámen certávi, cursum consummávi, fidem servávi. In réliquo repósita est mihi coróna justítiæ, quam reddet mihi Dóminus in illa die, justus judex: non solum autem mihi, sed et iis, qui díligunt advéntum ejus.}{english}{\noindent Beloved: I charge you, in the sight of God and Christ Jesus, who will judge the living and the dead by his coming and by his kingdom, preach the word, be urgent in season, out of season; reprove, entreat, rebuke with all patience and teaching. For there will come a time when they will not endure the sound doctrine; but having itching ears, will heap up to themselves teachers according to their own lusts, and they will turn away their hearing from the truth and turn aside rather to fables. But be watchful in all things, bear with tribulation patiently, work as a preacher of the Gospel, fulfill your ministry. As for me, I am already being poured out in sacrifice, and the time of my deliverance is at hand. I have fought the good fight, I have finished the course, I have kept the faith. For the rest, there is laid up for me a crown of justice, which the Lord, the just Judge, will give to me in that day; yet not to me only, but also to those who love his coming.}

\proper{Graduale \& Alleluia}{Ps. 36, 30 -- 31; Eccli. 45, 9}

\texteparacol{\lettrine{O}{s} justi meditábitur sapiéntiam, et lingua ejus loquétur judícium.
℣. Lex Dei ejus in corde i\-psíus: et non supplantabúntur gressus ejus.

\lettrine{A}{l}elúia, allelúia. Amávit eum Dóminus, et ornávit eum: stolam glóriæ índuit eum. Allelúia.}{english}{\noindent The mouth of the just shall meditate wisdom and his tongue shall speak judgment.
℣. The law of his God is in his heart, and his steps shall not be supplanted.

\noindent Alleluia, alleluia. The Lord loved him and beautified him and clothed him with a robe of glory. Alleluia.}

\reading{Evangelium}{Matt. 5, 3 -- 19}

\texteparacol{\lettrine{I}{n} illo témpore: Dixit Jesus discípulis suis: Vos estis sal terræ. Quod si sal evanúerit, in quo saliétur? Ad níhilum valet ultra, nisi ut mittátur foras, et conculcétur ab homínibus. Vos estis lux mundi. Non potest cívitas abscóndi supra montem pósita. Neque accéndunt lucérnam, et ponunt eam sub módio, sed super candelábrum, ut lúceat ó\-mnibus, qui in domo sunt. Sic luceat lux vestra coram homínibus, ut vídeant ópera vestra bona, et gloríficent Patrem vestrum, qui in cælis est. Nolíte putáre, quóniam veni sólvere legem aut prophétas: non veni sólvere, sed adimplére. Amen, quippe dico vobis, donec tránseat cælum et terra, jota unum aut unus apex non præteríbit a lege, donec ó\-mnia fiant. Qui ergo sólverit unum de mandátis istis mínimis, et docúerit sic hómines, mínimus vocábitur in regno cælórum: qui autem fécerit et docúerit, hic magnus vocábitur in regno cælórum.}{english}{\noindent At that time, Jesus said to his disciples: You are the salt of the earth. But if the salt lose its savour, wherewith shall it be salted? It is good for nothing any more but to be cast out, and to be trodden on by men. You are the light of the world. A city seated on a mountain cannot be hid. Neither do men light a candle and put it under a bushel, but upon a candlestick, that it may shine to all that are in the house. So let your light shine before men, that they may see your good works, and glorify your Father who is in heaven. Do not think that I am come to destroy the law, or the prophets. I am not come to destroy, but to fulfill. For amen I say unto you, till heaven and earth pass, one jot, or one tittle shall not pass of the law, till all be fulfilled. He therefore that shall break one of these least commandments, and shall so teach men, shall be called the least in the kingdom of heaven. But he that shall do and teach, he shall be called great in the kingdom of heaven.}

\proper{Offertorium}{Ps. 91, 13}

\texteparacol{\lettrine{J}{u}stus ut palma florébit: sicut cedrus, quæ in Líbano est, multiplicábitur}{english}{\noindent The just shall flourish like the palm tree: he shall grow up like the cedar of Libanus.}

\masspart{Secreta}
\texteparacol{\lettrine{S}{a}ncti Ephræm Confessóris tui atque Do\-ctóris nobis, Dómine, pia non desit orátio: quæ et múnera nostra concíliet; et tuam nobis indulgéntiam semper obtíneat.
Per Dóminum.

\lettrine{R}{e}spice, quǽsumus, Dómine, ad ineffábilem Cordis dilé\-cti Fílii tui caritátem: ut quod offérimus sit tibi munus accé\-ptum et nostrórum expiátio deli\-ctórum.

\lettrine[ante=2]{M}{u}nera tibi, Dómine, dicáta sanctífica: et, intercedéntibus san\-ctis Martýribus tuis Marco et Marcelliáno, per éadem nos placátus inténde.
Per Dóminum.}{english}{\noindent May the pious prayer of holy Ephrem thy confessor and illustrious Doctor be not wanting to us, o Lord, but make our offerings acceptable to thee and ever win for us thy mercy.
Through our Lord.

\noindent  2. Look, we beseech thee, o Lord, upon the heart of thy beloved Son, with its boundless love, so that what we offer, may be an acceptable gift in thy sight and an atonement for our sins.

\noindent 3. Sanctify the offerings dedicated to You, O Lord, and through the intercession of blessed Mark and Marcellianus, Your Martyrs, look upon us with mercy because of them.
Through our Lord.}

\masspart{Præfatio de sacratissimo Cordis Jesu}

\texteparacol{\lettrine{V}{e}re dignum et justum est, æquum et salutáre, nos tibi semper et ubíque grátias ágere: Dómine san\-cte, Pater omnípotens, ætérne Deus: Qui Unigénitum tuum, in Cruce pendéntem, láncea mílitis transfígi voluísti: ut apértum Cor, divínæ largitátis sacrárium, torréntes nobis fúnderet miseratiónis et grátiæ: et, quod amóre nostri flagráre numquam déstitit, piis esset réquies et pæniténtibus pater et salútis refúgium. Et ídeo cum Angelis et Archángelis, cum Thronis et Dominatiónibus cumque omni milítia cæléstis exércitus hymnum glóriæ tuæ cánimus, sine fine dicéntes:}{english}{\noindent It is truly meet and just, right and for our salvation, that we should at all times, and in all places, give thanks unto Thee, O holy Lord, Father almighty, everlasting God; Who didst will that Thine only-begotten Son, while hanging on the cross, should be pierced by a soldier's spear, that the Heart thus opened, a shrine of divine bounty, should pour out on us streams of mercy and grace, and that what never ceased to burn with love for us, should be a resting-place to the devout, and open as a refuge of salvation to the penitent. And therefore with angels and archangels, with thrones and dominions, and with all the hosts of the heavenly army, we sing the hymn of thy glory, evermore saying:}

%\rubrique{Præfatio, \normaltext{\pageref{præfatio}.}}

\proper{Communio}{Luc. 12, 42}

\texteparacol{\lettrine{F}{i}délis servus et prudens, quem constítuit dóminus super famíliam suam: ut det illis in témpore trítici mensúram.}{english}{\noindent The faithful and prudent servant whom the master will set over his household to give them their ration of grain in due time.}

\masspart{Postcommunio}
\texteparacol{\lettrine{U}{t} nobis, Dómine, tua sacrifícia dent salútem: beátus E\-phræm Conféssor tuus et Doctor egrégius, quǽsumus, precátor accédat.
Per Dóminum.

\lettrine[ante=2]{P}{r}ǽbeant nobis, Dómine Jesu, divínum tua san\-cta fervórem: quo dulcíssimi Cordis tui suavitáte percé\-pta; discámus terréna despícere, et amáre cæléstia:

\lettrine[ante=3]{S}{a}lutáris tui, Dómine, múnere satiáti, súpplices exorámus: ut, cujus lætámur gustu, intercedéntibus san\-ctis Martýribus tuis Marco et Marcelliáno, renovémur effé\-ctu.
Per Dóminum.}{english}{\noindent May blessed Ephrem thy Confessor and illustrious Doctor intercede for us, o Lord, that this thy sacrifice may obtain for us salvation. Through our Lord.

\noindent 2. May thy sacrament, o Lord Jesus, give us holy zeal, so that, seeing the sweetness of thy most loving heart, we may learn to despise the things of earth and love those of heaven.

\noindent 3. Refreshed by the sacramental gift of thy salvation, we humbly beseech thee, o Lord, through the intercession of blessed Mark and Marcellianus, thy martyrs, that, enjoying its savor, we may be made new by its action. Through our Lord.}

\office{ad vesperas}
\addcontentsline{toc}{subsection}{Ad  Vesperas} \phantomsection

\gscore[]{℣.}{Deus_In_Adjutorium_Festivus}
\zerobaroffsettextleft
\gscore[1. Ant.]{1. g}{an_suavi_jugo_solesmes_1961}
 %% watch in case initial remains low after a couple runs
 \resetbaroffsettextleft
\psalmus{109}{109_1g}{109_1}
 \smallscore{an_suavi_jugo}
 
 \gscore[2. Ant.]{2. D}{an_misericors_et_miserator_solesmes_1961}
\psalmus{110}{110_2}{110_2}
\smallscore{an_misericors_et_miserator}
 
  \gscore[3. Ant.]{7. a}{an_exortum_est_sacred_heart_solesmes_1961}
  \psalmus{111}{111_7a}{111_7}
  \smallscore{an_exortum_est_sacred_heart}
  
   \gscore[4 Ant.]{8. c}{an_quid_retribuam_solesmes_1961}
  \psalmus{115}{115_8c}{115_8}
  \smallscore{an_quid_retribuam}
   
    \gscore[5. Ant.]{4. A*}{an_apud_dominum_propitiatio_solesmes_1961}
    \psalmus{129}{129_4A_A_star}{129_4A_A_star}
    \smallscore{an_apud_dominum_propitiatio}
       
       \admagnificat
       
           \gscore[]{1. D}{an_ignem_solesmes_1961}
           
      \smallscore{Magnificat_1D}
      
         \magnificattexte{Magnificat_1_ordinaire}
         
    \smallscore{an_ignem}
    
      \vv Dóminus vobíscum. \rr Et cum spíritu tuo.
  
  \oratio
  
  Orémus.
  
  \lettrine{D}{e}us, qui nobis in Corde Fílii tui, nostris vulneráto peccátis, infinítos dile\-ctiónis thesáuros misericórditer largíri dignáris:~† concéde quǽsumus; ut illi devótum pietátis nostræ præstántes obséquium,~* dignæ quoque satisfa\-ctiónis exhibeámus offícium. Per eúmdem Dóminum.
  
  \rr Amen.
  
       \vv Dóminus vobíscum. \rr Et cum spíritu tuo.
       
       \gscore[Ant.]{2.}{an_o_doctor_optime_solesmes_1961_ephraem}
       
         \vv Justum dedúxit Dóminus per vias rectas.
  
\rr Et osténdit illi regnum Dei.

 Orémus.

\lettrine{D}{e}us, qui Ecclésiam tuam beáti Ephræm Confessóris tui et Do\-ctóris mira eruditióne et præcláris vitæ méritis illustráre voluísti:~† te súpplices exorámus; ut i\-pso intercedénte,~* eam advérsus erróris et pravitátis insídias perénni tua virtúte deféndas.

\gscore[Ant.]{8.}{an_cs_veni_sponsa_i_vesp_solesmes_1961}

\vv Spécie tua et pulchritúdine tua.

\rr Inténde, próspere procéde, et regna.
\newpage
Orémus.

\lettrine{D}{e}us, qui beátam Juliánam Vírginem tuam extrémo morbo laborántem, pretióso Fílii tui Córpore mirabíliter recreáre dignátus es:~† concéde, quǽsumus; ut ejus intercedéntibus méritis, nos quoque eódem in mortis agóne refé\-cti ac roboráti,~* ad cæléstem pátriam perducámur.
   
    \gscore[Ant.]{8.}{an_cs_istorum_est_enim_solesmes_1961}
    
  \vv  Lætámini in Dómino, et exsultáte justi.
    
\rr Et gloriámini omnes recti corde.

Orémus.

\lettrine{D}{e}us, qui nos ánnua san\-ctórum Mártyrum túorum Gervásii et Protásii sole\-mnitáte lætíficas:~† concéde propítius; ut quorum gaudémus méritis,~* accendámur exémplis. Per Dóminum.

\rr Amen.
    
\gscore[]{2.}{BD_Duplex_I_Vesp}
   
  \vv Fidélium ánimæ per misericórdiam Dei requiéscant in pace.
  
  \rr Amen.
  
 \rubrique{Tunc dicitur Completorium. Si non dicendum:}
  
  \secreto
  
\vv Dóminus det nobis suam pacem.
  
\rr Et vitam ætérnam. Amen.

\rubrique{\normaltext{Salve Regína,} etc., \normaltext{\pageref{an_salve_regina_solesmes}.}}
